
\subsection{Green 张量统一前向模型}\label{subsec:green-forward-closure-dp47}

线性介质对任意电流源的响应已经由推迟 Green 张量确定。自由电子实验只是改变了源和读出方式：电子自身作为源时，可以统计它失去的总能量，也可以只统计传到远场的辐射；外加相干光场作为源时，则读取电子获得的相干能量调制。因此三类观测共享同一个频域响应核，只在源项和探测投影上不同。对给定经典电子轨迹 $\mathbf r_e(t)$，电流 Fourier 分量为
\begin{equation}
 \widetilde{\mathbf J}_e(\mathbf r,\omega)
 =q_{\mathrm e}\int\!\dd t\,
 \mathbf v(t)\delta^{(3)}[\mathbf r-\mathbf r_e(t)]\ee^{\ii\omega t}.
\label{eq:electron_current_fourier_dp47}
\end{equation}
由线性响应产生的诱导场为
\begin{equation}
\boxed{
 \widetilde{\mathbf E}_{\rm ind}(\mathbf r,\omega)
 =\ii\mu_0\omega
 \int\!\dd^3 r'\,
 \mathbf G_{\rm sc}^R(\mathbf r,\mathbf r';\omega)
 \widetilde{\mathbf J}_e(\mathbf r',\omega),}
\label{eq:induced_field_green_dp47}
\end{equation}
其中 $\mathbf G_{\rm sc}^R=\mathbf G^R-\mathbf G_0^R$ 只保留环境相对于参考真空的散射响应。

电子单位角频率损失概率可由电子对诱导场所做功得到。对匀速直线轨迹 $\mathbf r_e(t)=\mathbf R_0+\mathbf v_0t$，写成 Green 张量双时间积分为
\begin{equation}
\boxed{
 \Gamma_{\rm EELS}(\omega)
 =\frac{\mu_0q_{\mathrm e}^2}{\pi\hbar}
 \int\!\dd t\,\dd t'\,
 \ee^{\ii\omega(t-t')}
 \mathbf v_0\!\cdot
 \operatorname{Im}\mathbf G_{\rm sc}^R
 (\mathbf r_e(t),\mathbf r_e(t');\omega)
 \cdot\mathbf v_0 .}
\label{eq:eels_green_master_dp47}
\end{equation}
$\Gamma_{\rm EELS}(\omega)\dd\omega$ 是无量纲概率，因此 $[\Gamma_{\rm EELS}]={\rm s}$。

CL 不是把\eqnrefs{eq:eels_green_master_dp47} 中的 $\operatorname{Im}\mathbf G^R$ 换成模方，而是先把电子电流经 Green 张量传播到远场探测通道。令探测模式 $\lambda$ 的归一化接收函数为 $\mathbf f_\lambda(\mathbf r,\omega)$，定义通道振幅
\begin{equation}
 \mathcal A_\lambda(\omega)
 =\ii\mu_0\omega
 \int\!\dd^3 r\,\dd^3 r'\,
 \mathbf f_\lambda^*(\mathbf r,\omega)\!\cdot
 \mathbf G^R(\mathbf r,\mathbf r';\omega)
 \widetilde{\mathbf J}_e(\mathbf r',\omega).
\label{eq:cl_channel_amplitude_dp47}
\end{equation}
在接收模式按单光子能量归一化后，谱计数率具有
\begin{equation}
\boxed{
 \Gamma_{\rm CL}^{\rm det}(\omega)
 =\sum_{\lambda\in\mathcal C}
 \frac{|\mathcal A_\lambda(\omega)|^2}{\hbar\omega},}
\label{eq:cl_detected_green_master_dp47}
\end{equation}
其中集合 $\mathcal C$ 编码数值孔径、偏振、频率窗和空间模式匹配。这样“总辐射”与“被探测 CL”严格区分。

受激 PINEM 则由外加相干场沿电子轨迹所做的功决定。对沿 $z$ 方向以速度 $v_0$ 传播的电子，时域功为
\begin{equation}
\boxed{
 \Delta E_{\rm ext}
 =q_{\mathrm e}\int_{-\infty}^{\infty}\!\dd t\,
 v_0 E_{{\rm ext},z}(\mathbf R_0,v_0t;t).}
\label{eq:pinem_time_domain_work_dp68}
\end{equation}
若外场只含角频率 $\omega$ 的单频分量，并以 $\mathbf E_{{\rm ext},\omega}(\mathbf r)$ 表示其复幅，那么沿轨迹抽取这一频率分量后，经典无量纲耦合为
\begin{equation}
\boxed{
 \beta(\omega)
 =\frac{q_{\mathrm e}}{\hbar\omega}
 \int_{-\infty}^{\infty}\!\dd z\,
 E_{{\rm ext},\omega,z}(\mathbf R_0,z)
 \ee^{-\ii\omega z/v_0}.}
\label{eq:pinem_beta_green_unified_dp47}
\end{equation}
若外场本身由 Fourier 源电流 $\widetilde{\mathbf J}_{\rm ext}$ 经同一环境产生，则将
\begin{equation*}
 \mathbf E_{{\rm ext},\omega}
 =\ii\mu_0\omega\,\mathbf G^R*\widetilde{\mathbf J}_{\rm ext}
\end{equation*}
代入\eqnrefs{eq:pinem_beta_green_unified_dp47}，PINEM 也成为同一个 $\mathbf G^R$ 的线性泛函。

因此三类实验的结构可压缩为
\begin{equation}
\boxed{
 \mathbf G^R
 \longrightarrow
 \begin{cases}
 \operatorname{Im}\mathbf G^R & \text{EELS：电子不可逆损失},\\
 \Pi_{\mathcal C}\mathbf G^R\widetilde{\mathbf J}_e & \text{CL：辐射到收集通道},\\
 \mathbf G^R\widetilde{\mathbf J}_{\rm ext} & \text{PINEM：外场相干驱动},
 \end{cases}}
\label{eq:green_three_experiments_map_dp47}
\end{equation}
这里的差别来自输入源和最终读取的观测量，而不是三种不同的介质理论。

实际反演还要把理论谱映射到探测器计数。把待估参数记为 $\boldsymbol\vartheta$，理论连续谱记为 $s(\omega|\boldsymbol\vartheta)$，仪器响应矩阵为 $R_{j}(\omega)$，则第 $j$ 个像素的期望计数为
\begin{equation}
\boxed{
 \lambda_j(\boldsymbol\vartheta)
 =N_e\int\!\dd\omega\,
 R_j(\omega)s(\omega|\boldsymbol\vartheta)+b_j,}
\label{eq:detector_count_forward_dp47}
\end{equation}
其中 $N_e$ 是入射电子数，$b_j$ 是背景期望计数。若计数服从独立 Poisson 统计，完整似然为
\begin{equation}
\boxed{
 \log\mathcal L(\boldsymbol\vartheta)
 =\sum_j\left[n_j\log\lambda_j-\lambda_j-\log(n_j!)\right].}
\label{eq:poisson_likelihood_green_inverse_dp47}
\end{equation}
这使 Green 张量、电子态、收集模式、仪器分辨率与最终参数反演处在同一个前向模型中。


\begin{derivation}[从\eqnrefs{eq:pinem_time_domain_work_dp68}到\eqnrefs{eq:pinem_beta_green_unified_dp47}]
采用全文的正频率复幅约定，把真实外场写成
\[
E_{{\rm ext},z}(\mathbf r,t)
=\int_0^\infty\frac{\dd\omega}{2\pi}
\left[
E_{{\rm ext},\omega,z}(\mathbf r)\ee^{-\ii\omega t}
+E_{{\rm ext},\omega,z}^*(\mathbf r)\ee^{+\ii\omega t}
\right].
\]
沿匀速轨迹 $\mathbf r_e(t)=(\mathbf R_0,v_0t)$ 代入正文\eqnrefs{eq:pinem_time_domain_work_dp68}。对固定正频率 $\omega$，定义复功幅
\begin{align*}
\mathcal W_\omega
&\equiv
q_{\mathrm e}\int_{-\infty}^{\infty}\dd t\,
 v_0 E_{{\rm ext},\omega,z}(\mathbf R_0,v_0t)
 \ee^{-\ii\omega t}\\
&=q_{\mathrm e}\int_{-\infty}^{\infty}\dd z\,
 E_{{\rm ext},\omega,z}(\mathbf R_0,z)
 \ee^{-\ii\omega z/v_0},
\end{align*}
其中第二行只用了 $z=v_0t$ 与 $\dd z=v_0\dd t$。指数因子表明电子只抽取外场纵向 Fourier 谱在
\[
k_z=\frac{\omega}{v_0}
\]
这一同步线上对应的分量。

单个边带步长的能量是 $\hbar\omega$，因此把相干功幅除以这一能量定义无量纲生成元
\[
\beta(\omega)\equiv\frac{\mathcal W_\omega}{\hbar\omega},
\]
立刻得到正文\eqnrefs{eq:pinem_beta_green_unified_dp47}。这个除法不是经验归一化：在能量格演化算符中，$\beta$ 必须是乘在无量纲平移生成元前的无量纲系数。

若外场由外源电流产生，正文采用
\[
\mathbf E_{{\rm ext},\omega}(\mathbf r)
=\ii\mu_0\omega
\int\dd^3r'\,
\mathbf G^R(\mathbf r,\mathbf r';\omega)
\cdot\widetilde{\mathbf J}_{\rm ext}(\mathbf r',\omega).
\]
把它代回 $\beta$ 并约去 $\omega$，得到
\begin{align*}
\beta(\omega)
={}&\frac{\ii\mu_0q_{\mathrm e}}{\hbar}
\int\dd z\dd^3r'\,
 \ee^{-\ii\omega z/v_0}
\hat{\mathbf z}\cdot
\mathbf G^R((\mathbf R_0,z),\mathbf r';\omega)
\cdot\widetilde{\mathbf J}_{\rm ext}(\mathbf r',\omega).
\end{align*}
因此 PINEM 耦合与 EELS/CL 使用的是同一个推迟 Green 算符；PINEM 由外源产生相干场并沿电子轨迹读取线性振幅，EELS 则对电子自身诱导场取耗散二次型。
\end{derivation}


一般的“做功 $\to\operatorname{Im}\mathbf G^R\to$ 损失概率”关系已经由\eqnrefs{eq:prescribed_current_green_loss_master,eq:green_tensor_eels_master} 建立。对具体几何，只需把相应响应核代入一般表达式：使用散射 Green 张量 $\mathbf G_{\rm sc}^R=\mathbf G^R-\mathbf G_0^R$ 去除参考真空贡献，并取匀速轨迹
\[
\mathbf r_e(t)=\mathbf R_0+\mathbf v_0t,
\qquad
\mathbf J_e(\mathbf r,t)=q_{\mathrm e}\mathbf v_0\delta^{(3)}[\mathbf r-\mathbf r_e(t)].
\]
一般表达式随即化成\eqnrefs{eq:eels_green_master_dp47}。由此可以直接看到：材料和几何只改变 Green 前向核，电子跃迁理论仍由同一个电流--响应双线性控制。

